\scanpage{55}

% 1–2
\[
 F(\mathbb E^i_{\bullet})\longrightarrow\sigma(\rho(F))(\mathbb E^i_{\bullet})=\operatorname{Hom}^{\bullet}_{k}(\mathbb E^{\deleted{\unsure{i}}\bullet}_{i},\underset{\substack{\Vert\\F(\mathbb E_{\bullet}^{\bullet})}}{\rho(F)})\simeq F(\mathbb E^i_{\bullet})
\]

% 3–6
\[
 \begin{aligned}
 F(\Psi^{i}_{\bullet})&\longrightarrow\operatorname{Hom}^{\bullet}_{k}(\Psi^{\deleted{\unsure{i}}\bullet}_{i},F(\mathbb E^{\bullet}))\simeq Z^{i}(F(\mathbb E^{\bullet}))\\
 &\simeq\operatorname{Ker}(F(\mathbb E^i)\to F(\mathbb E^{i+1}))\\
 &\simeq F(\underbrace{\operatorname{Ker}(\mathbb E^i\longrightarrow\mathbb E^{i+1})}_{\Sigma\text{-}\operatorname{Ker}\,!})=F(\Psi^i)
 \end{aligned}
\]

% 7–13
Donc on voit que $u(L_{\bullet}):F(L_{\bullet})\to\sigma\rho(F)(L_{\bullet})$ est un iso sur les $\mathbb E^i$ et les $\Psi^{i}_{\bullet}$, donc sur \deleted{la} la cat. additive $\mathbb P_0\supset\mathbb K$ des sommes de tels.) Pour avoir $\forall L_{\bullet}$, on \inserted{essaye en} raisonnant par récurrence sur la longueur $n$ \deleted{\unsure{de $L_{\bullet}$|d'un $L_{\bullet}$}}, \inserted{de l'intervalle $[p,q]$ tel que $L^i=0$ si $i\notin[p,q]$.} C'est vu si $n=0$, ($L_{\bullet}$ somme \deleted{\unsure{de $\Psi_{\bullet}$}} \inserted{$\simeq(\Psi^{i}_{\bullet})^{r}$}), on suppose $n>0$ et supposons que c'est fait pour $n-1$.

% 14–17, truncation diagram; brace under the truncated terms
\diagram{\begin{tikzcd}[ampersand replacement=\&,column sep=.55em,execute at end picture={\path (trleft.south west) -- (trright.south east) coordinate[midway] (trmiddle); \draw (trleft.south west) .. controls +(0,-.4em) and +(-.8em,.3em) .. ([yshift=-.65em]trmiddle) node[below] {$\text{tronqué }L'_{\bullet}$} .. controls +(.8em,.3em) and +(0,-.4em) .. (trright.south east);}]
 \&\&\&\cdots\arrow[r]\&0\arrow[r]\&L_p\arrow[r,"\mathrm{id}"]\&L_p\arrow[r]\&0\&\& \left.\vphantom{\begin{matrix}x\\x\end{matrix}}\right\}\ (\mathbb E_{\bullet}^{p\deleted{\unsure{+1}}})^{r}\deleted{\unsure{L_{\bullet}''}}\\
 \cdots\arrow[r]\&0\arrow[r]\&|[alias=trleft]|L_{p+n}\arrow[r]\&L_{p+n-1}\arrow[r]\&\cdots\arrow[r]\&|[alias=trright]|L_{p+1}\arrow[u,"{d_{p+1}}"]\arrow[r]\&L_p\arrow[u,"{\mathrm{id}_{L_p}}"]\arrow[r]\&\deleted{\unsure{0}}\arrow[r]\&0\& L_{\bullet}\arrow[u]
\end{tikzcd}}

% 18
on trouve un \unsure{mor|m} de suites exactes (où $L_p\simeq k^r$)

% 19–20, square with exact rows
\diagram{\begin{tikzcd}[ampersand replacement=\&,column sep=1.2em]
 0\arrow[r]\&(\Psi^p_{\bullet})^r\arrow[r]\arrow[d,equal]\&L_{\bullet}\arrow[r]\arrow[d]\&\deleted{\unsure{\Psi}}L_{\bullet}^{\unsure{\prime|\prime\prime}}\arrow[r]\arrow[d]\&0\\
 0\arrow[r]\&(\Psi^p_{\bullet})^r\arrow[r]\&(\mathbb E_{\bullet}^{p})^r\arrow[r]\&(\Psi_{\bullet}^{p+1})^r\arrow[r]\&0
\end{tikzcd}}

% 19–24, starts to right of square
qui exprime $L_{\bullet}$ comme produit fibré de $L'_{\bullet}$ et $(\mathbb E_{\bullet}^{p})^r$ sur $(\Psi_{\bullet}^{p+1})^r$ (i.e. comme noyau d'un hom \deleted{surj.} $\Sigma$-\uline{épi} $(\mathbb E_{\bullet}^{p})^{r}\times L_{\bullet}^{\unsure{1|\prime}}\longrightarrow(\Psi_{\bullet}^{p+1})^{r}$). Donc on conclut par

% archive, bottom left
% archive sheet number: 54
